\documentclass[10pt,a4paper]{amsart}
\usepackage[margin=19mm]{geometry}
\usepackage{amsmath,amssymb,mathtools}
\usepackage[scr=rsfs,cal=euler]{mathalfa}
\usepackage{xfrac}
\usepackage[unicode,hidelinks,bookmarks=false]{hyperref}
\newcommand{\ie}{i.e.\ }
\newcommand{\cf}{cf.\ }
\newcommand{\resp}{respectively\ }
\newcommand{\Subsection}{\subsection}
\providecommand{\nobreakdash}{\nobreak\hbox{-}}
\setlength{\emergencystretch}{2em}
\multlinegap0pt
\usepackage{amssymb}
\usepackage{mathtools}
\usepackage[shortlabels]{enumitem}
\usepackage{tikz}
\newtheorem{assumption}{Assumption}
\newtheorem{definition}{Definition}
\newtheorem{theorem}{Theorem}
\newcommand{\Chat}{\hat{C}}
\newcommand{\E}{\mathbb{E}}
\newcommand{\F}{\mathcal{F}}
\newcommand{\Fhat}{\hat{\F}}
\newcommand{\KL}[2]{D_{\mathrm{KL}}(#1 \,\|\, #2)}
\title{\bfseries Polynomial Context-Truncation Sensitivity in Autoregressive Language Models:\\[2pt] Sequential Wyner--Ziv Bounds for KV Cache Compression}
\author{Munsik Kim}
\date{Source: arXiv:2605.25085v2 (CC BY 4.0)}
\begin{document}
\maketitle

\noindent\emph{Source and licence.} \bfseries Polynomial Context-Truncation Sensitivity in Autoregressive Language Models:\\[2pt] Sequential Wyner--Ziv Bounds for KV Cache Compression. Version arXiv:2605.25085v2 (CC BY 4.0), distributed by arXiv under the Creative Commons Attribution 4.0 International licence, http://creativecommons.org/licenses/by/4.0/, as recorded in the arXiv licence metadata for this exact version and shown on the abstract page https://arxiv.org/abs/2605.25085v2. Full licence notice: \texttt{licenses/arxiv-2605.25085-CC-BY-4.0.txt}.\par\smallskip

\noindent\textit{Editorial scope.} Counted unit: the article's theorem thm:asymptotic (source0.tex lines 568--574). The article's complete original proof of it is delivered with it (source0.tex lines 2638--2684, one \texttt{proof} environment of 47 lines, which the article places away from the statement). No mathematics is written, repaired or re-derived: the statement and the proof are the article's own lines, in the article's own notation, and no retained line is substituted or re-flowed.  Retained uncounted as the source-local context the proof needs: lines 417--421; lines 465--481.  Nothing was omitted from any retained span, so the package carries no omission record.  No retained line contains a citation, so the wrapper carries no bibliography and no bibliography entry.  No retained line contains a figure environment, an image-inclusion command or drawing code, so no image is delivered and the package declares no asset dependency.  Editorial formatting only, in addition: an AMS \texttt{amsart} preamble that restates the article's own declarations and macros used by the retained lines, namely the environments \texttt{assumption}, \texttt{definition}, \texttt{theorem} on separate counters, together with the standard packages those lines need.  The retained environments print in this document as Assumption 1, Definition 1, Theorem 1 in order of appearance, whereas the article prints the same items with the numbering of its own full document.  The complete native source of the article as distributed in the arXiv e-print for this exact version is bundled unchanged as \texttt{sources/201323/source0.tex}.
\begin{assumption}[Bounded logits]\label{ass:bounded-logits}
There exists $B(n_{\max}) > 0$ such that
$\|Z_t\|_\infty \leq B(n_{\max})$ a.s.\ for $t \leq n_{\max}$;
consequently, $p_t \geq \epsilon := V^{-1}e^{-2B}$ componentwise.
\end{assumption}

\begin{definition}[Sequential WZ rate function]
\label{def:Rstar}
An \emph{auxiliary process} $\{U_t\}$ is an $\F_t$-adapted sequence
that is $\Fhat_t$-measurable. The per-step pointwise WZ rate is
$\mathcal{R}_t^{\mathrm{WZ}}
:= I(X_t; U_t \mid \Fhat_{t-1}) - I(U_t; Q_t \mid \Fhat_{t-1})$.
The auxiliary is \emph{$D$-admissible} if some measurable decoder
$\psi_t = g_t(U_t, Q_t, \Fhat_{t-1})$ achieves
$\E[\KL{p_t}{\tilde p_t}] \leq D$ on average (the decoder may use the
realized code history, which is $\Fhat_{t-1}$-measurable). The sequential
Wyner--Ziv rate function is
\[
R^\star(D)
:= \inf_{\{U_t\}\ D\text{-admissible}}
\liminf_{n} \frac{1}{n}\sum_{t=1}^n \E[\mathcal{R}_t^{\mathrm{WZ}}].
\]
\end{definition}

\begin{theorem}[Asymptotic lower bound]\label{thm:asymptotic}
Under Assumption~\ref{ass:bounded-logits}, for any
encoder--decoder pair with $\limsup_n \E[D_n] \leq D$,
\[
\liminf_{n} \E[R_n] \;\geq\; R^\star(D).
\]
\end{theorem}

\begin{proof}[Proof of Theorem~\ref{thm:asymptotic}]
\textbf{Step 1 (rate $\to$ MI by chain rule).}
For each $t$,
\[
H(\Chat_t \mid \Fhat_{t-1})
\geq I(X_{\leq t}; \Chat_t \mid \Fhat_{t-1})
\]
by the standard inequality $H(\Chat_t \mid \Fhat_{t-1})
= I(\Chat_t; X_{\leq t} \mid \Fhat_{t-1})
+ H(\Chat_t \mid X_{\leq t}, \Fhat_{t-1})
\geq I(X_{\leq t}; \Chat_t \mid \Fhat_{t-1})$.

\textbf{Step 2 (single-letterization via chain rule).}
With $U_t := \Chat_t$,
\[
\begin{aligned}
I(X_{\leq t}; \Chat_t \mid \Fhat_{t-1})
&= I(X_t; \Chat_t \mid \Fhat_{t-1}) \\
&\quad + I(X_{<t}; \Chat_t \mid X_t, \Fhat_{t-1}) \\
&\geq I(X_t; \Chat_t \mid \Fhat_{t-1}),
\end{aligned}
\]
by non-negativity of the second term. Adding and subtracting
$\mathcal{I}_t^{UQ}$:
\[
I(X_t; \Chat_t \mid \Fhat_{t-1})
= \mathcal{R}_t^{\mathrm{WZ}} + \mathcal{I}_t^{UQ} \geq \mathcal{R}_t^{\mathrm{WZ}},
\]
using $\mathcal{I}_t^{UQ} \geq 0$.

\textbf{Step 3 (admissibility of $\{\Chat_t\}$).}
Under the hypothesis $\limsup_n \E[D_n] \leq D$, the encoder--decoder
pair with $U_t = \Chat_t$ and decoder
$g_t(\Chat_t, Q_t, \Fhat_{t-1}) = \psi_t(\Chat_{\le t}, Q_t)$---valid
since the past reconstructions $\Chat_{<t}$ are
$\Fhat_{t-1}$-measurable---is $D$-admissible per
Definition~\ref{def:Rstar}.

\textbf{Step 4 (infimum gives $R^\star(D)$).}
Taking expectations and averaging,
\[
\E[R_n] = \tfrac{1}{n}\sum_t \E[H(\Chat_t \mid \Fhat_{t-1})]
\geq \tfrac{1}{n}\sum_t \E[\mathcal{R}_t^{\mathrm{WZ}}].
\]
By Definition~\ref{def:Rstar} and Step~3,
$\liminf_n \tfrac{1}{n}\sum_t \E[\mathcal{R}_t^{\mathrm{WZ}}] \geq R^\star(D)$.
\end{proof}

\end{document}
